\section{Introduction}
Advanced Persistent Threats (APTs) are stealthy, multi-stage cyber campaigns that compromise enterprise systems over extended periods. Because APTs execute benign-looking commands and proceed through multi-hop lateral movements, traditional signature-based security systems are ineffective. Consequently, host audit logging and provenance graph analysis have emerged as powerful paradigms for APT detection. A provenance graph models system events (e.g., file reads, process spawns, network flows) as directed edges, preserving the causal history of system execution.

Despite their potential, existing state-of-the-art provenance-based IDS---such as OCR-APT \cite{ocrapt2025}, TraceCluster \cite{tracecluster2026}, and StageFinder \cite{stagefinder2026}---share a fatal limitation: \textbf{the clean-data assumption}. They require a pristine ``benign'' period or fully labeled training dataset to learn a statistical baseline of normal behavior. This assumption is unrealistic in practice because: (1) attackers may already be active inside the environment (stealthy contamination), (2) systems evolve naturally (concept drift), (3) logs are too massive to label manually, and (4) static anomaly scoring triggers false positive fatigue.

\textbf{What this paper makes possible.}~Prior to this work, no provenance-based IDS could achieve competitive detection on real enterprise logs (e.g., OpTC AUC $>$0.95) without either a clean baseline collection period or fully labeled training data. Our framework demonstrates that streaming provenance telemetry can be projected into a formal causal behavioral state space using online SCM discovery on raw, potentially contaminated streams, and that multi-perspective reasoning on this state space yields reliable APT detection with distribution-free false-alarm guarantees under concept drift. This eliminates the need for a clean commissioning window, enabling deployment from day one on live enterprise logs.

To address these challenges, we propose a \textbf{Behavioral State Inference} framework that infers latent cyber risk from evolving behavioral states on uncurated streams. Our core insight is that \textbf{causal relationships are more stable than statistical correlations}. However, pure causal anomaly detectors (such as a pure SCM model or Causal-IDS) evaluate causal mechanism violations in isolation (e.g., via simple p-values or regression residuals). This pure causal mapping is highly sensitive to statistical noise and completely ignores temporal patterns (e.g., slow-burn APT progression) and structural combinations.

To overcome this, we design a framework that projects system behavior into a formal \textbf{Behavioral State Space} $S_t = \phi(X_t, G_t) \in \mathcal{S}$ extracted from online Structural Causal Models (SCMs) learned directly from raw, unlabeled system logs using robust partial correlation. While the behavioral state coordinates are projected in a strictly unsupervised manner, system trajectories are analyzed in parallel by three reasoning operators:
\begin{enumerate}
\item \textbf{Local State Reasoning ($R_L$)}: Measures instantaneous mechanism violations on the state representation in an unsupervised manner.
\item \textbf{Temporal State Reasoning ($R_T$)}: Evaluates state trajectories over a sequence window using an unsupervised, lightweight LSTM Autoencoder to isolate slow, multi-stage temporal progression.
\item \textbf{Behavioral State Reasoning ($R_B$)}: Classifies complex, non-linear attack manifolds using a \textbf{Behavioral Pattern Reasoner (BPR)} (implemented via a Random Forest) calibrated in a semi-supervised manner on a small validation partition containing sparse threat labels.
\end{enumerate}

The outputs of these reasoning operators are combined in a \textbf{calibrated risk inference layer} using a stacked meta-fusion operator whose weights are trained on the same sparse validation partition. Final alerts are then calibrated using online conformal prediction to satisfy a user-defined False Positive Rate (FPR) budget.

\subsection{Novel Components Introduced}
\label{sec:novel_components}
Several components in our pipeline are standard and reused from prior work. For transparency, we explicitly distinguish novel from reused components:

\begin{itemize}
\item \textbf{Reused (prior art)}: PCMCI causal discovery algorithm \cite{pcmci2019}; RobustParCorr robust partial correlation (Tigramite library); linear structural causal models (standard econometrics); LSTM Autoencoder architecture (standard deep learning); Random Forest classifier (scikit-learn); logistic regression meta-learner; conformal prediction theory \cite{angelopoulos2021gentle}.
\item \textbf{Novel (this paper)}: (1) The \emph{12-dimensional Causal Behavioral State Space} $\mathcal{S}$ defined over five coordinate groups (Structural Consistency, Activity Dynamics, Behavioral Complexity, Mechanism Violation, Intervention Locality)---a formal state representation specifically designed for streaming provenance telemetry. (2) The \emph{Causal Intervention Score (CIS)}, a graph-boundary metric rooted in spectral graph theory \cite{chung1997spectral} that quantifies multi-hop violation propagation to separate administrative noise from APT lateral movement. (3) The \emph{Multi-Perspective Parallel Reasoning} architecture that combines local, temporal, and behavioral operators on the shared state space, resolving model-mismatch failures of pure causal detectors. (4) The \emph{online conformal adaptation with change-point-triggered queue flushing} that maintains distribution-free FPR bounds under concept drift.
\end{itemize}

By shifting to multi-perspective state reasoning, our framework leverages the best of both worlds: the robust, drift-invariant foundation of causal structure learning, and the high-discrimination capability of temporal and pattern-matching classifiers. This state reasoning approach yields consistent and significant performance improvements. For instance, on StreamSpot, where pure causal models fail completely due to SCM mismatch on heterogeneous graph databases (AUC 0.4991), our framework achieves AUC of 0.6909. On real enterprise Windows logs (OpTC), it reaches AUC 0.9628, surpassing Isolation Forest (0.5968) and deep Autoencoder (0.9364) baselines.

\textbf{Comparison with Prior Causal IDSs.}~Our work is closely related to two recent causal network IDS frameworks: \textbf{Causal-IDS} (2026)~\cite{causalids2026} and \textbf{CausalGraph} (2026)~\cite{causalgraph2026}.
Compared to Causal-IDS, which models network flow statistics using a static SCM fitted on fully benign data, our framework (1) executes online on uncurated streams without assuming clean baseline data, (2) recovers from concept drift via online SCM refitting, and (3) extracts a behavioral state representation that leverages non-linear classifiers rather than relying purely on noisy individual causal residuals.
Compared to CausalGraph, which utilizes LLMs to parse CTI and suggest counterfactual interventions for explanation, our model is designed for high-velocity streaming environments; it replaces expensive LLM-based prior/intervention generation with a lightweight, purely structural Causal Intervention Score (CIS) computed in milliseconds, achieving real-time operational feasibility while providing equivalent topological anomaly reasoning.


\subsection{Summary of Contributions}
\begin{itemize}
\item \textbf{Causal Behavioral State Representation}: We introduce a formal behavioral state representation derived from online causal invariants. Telemetry and SCM residuals are projected into a 12-dimensional state space representing structural consistency, mechanism violation, intervention locality, behavioral complexity, and activity dynamics.
\item \textbf{Novel CIS Metric}: We design the Causal Intervention Score, a graph-boundary-based metric that provides principled topological separation between localized administrative events and multi-hop APT propagation.
\item \textbf{Multi-Perspective State Reasoning}: We design a multi-perspective state reasoning framework that analyzes local, temporal, and behavioral manifestations of the same state in parallel, resolving the model-mismatch issues of pure causal anomaly detectors.
\item \textbf{Calibrated Risk Inference Layer}: We establish a calibrated risk inference layer utilizing stacked meta-fusion and online conformal calibration to provide bounded false alarm guarantees under concept drift.
\item \textbf{Operational Verification}: We benchmark our model on five diverse datasets (OpTC, BETH, StreamSpot, TC3, and NODLINK). Our framework achieves AUC of \textbf{0.9628} on OpTC, \textbf{0.9901} on BETH, and \textbf{0.9837} and \textbf{0.9784} on TC3 and NODLINK, outperforming all static unsupervised alternatives.
\item \textbf{Optimized Execution}: We implement vectorized CDF operations and binary search conformal checks, yielding a \textbf{13$\times$ speedup} and ensuring runtime readiness for high-velocity enterprise streams.
\end{itemize}
